\documentclass[12pt,a4paper]{article}

% Кодировка и язык
\usepackage[utf8]{inputenc}
\usepackage[T2A]{fontenc}
\usepackage[english,russian]{babel}

% Математические пакеты
\usepackage{amsmath}
\usepackage{amssymb}
\usepackage{amsfonts}
\usepackage{mathtools}

% Графика и таблицы
\usepackage{graphicx}
\usepackage{booktabs}
\usepackage{array}
\usepackage{multirow}
\usepackage{tabularx}

% Оформление
\usepackage{geometry}
\usepackage{titlesec}
\usepackage{fancyhdr}
\usepackage{hyperref}
\usepackage{tocloft}
\usepackage{enumitem}

% Поля страницы
\geometry{left=2.5cm,right=2.5cm,top=2.5cm,bottom=2.5cm}

% Нумерация страниц
\pagestyle{fancy}
\fancyhf{}
\fancyhead[LE,RO]{\thepage}
\fancyhead[LO]{Классическая модель спина электрона}
\fancyhead[RE]{Технический отчёт}

% Гиперссылки
\hypersetup{
    colorlinks=true,
    linkcolor=blue,
    filecolor=magenta,
    urlcolor=cyan,
    citecolor=blue,
}

% Заголовок документа
\title{\textbf{КЛАССИЧЕСКАЯ МОДЕЛЬ СПИНА ЭЛЕКТРОНА}\\[0.5cm]
\large На основе вращающегося скалярного потенциала\\[0.3cm]
\normalsize Интеграция с Динамической Теорией Электромагнитного Поля Максвелла (ДТПМ)}

\author{На основе дискуссии по моделированию\\
с участием теории С.А. Цикры}

\date{Март 2026 г.}

\begin{document}

\maketitle

\begin{abstract}
\noindent В данном отчёте представлена непротиворечивая классическая модель спина электрона, основанная на вращении поля скалярного потенциала вместе с зарядом. Показано, что правильное гиромагнитное отношение $g=2$ достигается при введении физического радиуса перехода $R_{tr}$, совпадающего с комптоновской длиной волны $\lambda_c$. Доказана невозможность интегрирования до бесконечности из-за логарифмической расходимости момента импульса. Проведён анализ теории С.А. Цикры и предложена схема интеграции двух подходов. Модель воспроизводит правильную дипольную асимптотику магнитного поля и объясняет квантоподобные эффекты (спиновые волны, обменное взаимодействие) в рамках классической электродинамики.

\vspace{0.3cm}
\noindent\textbf{Ключевые слова:} спин электрона, гиромагнитное отношение, вращающийся скалярный потенциал, комптоновская длина волны, ДТПМ, классическая электродинамика.
\end{abstract}

\tableofcontents
\newpage

%================================================================================
\section{Введение}
%================================================================================

В рамках классической электродинамики существует давняя проблема объяснения спина электрона и его гиромагнитного отношения $g \approx 2$. Традиционные подходы, основанные на модели вращающегося заряженного шара, сталкиваются с фундаментальными трудностями:

\begin{itemize}
    \item Необходимость сверхсветовых скоростей на поверхности электрона
    \item Расхождение теоретического значения $g$-фактора с экспериментальным
    \item Проблема $4/3$ при вычислении электромагнитной массы
\end{itemize}

Данная работа предлагает альтернативный подход, основанный на следующих ключевых идеях:

\begin{enumerate}
    \item \textbf{Вращение поля скалярного потенциала} вместе с зарядом, а не только вращение самого заряда
    \item \textbf{Использование формулы импульса поля по Рорлиху} вместо стандартного вектора Пойнтинга
    \item \textbf{Введение физического радиуса перехода} $R_{tr}$, на котором линейная скорость поля асимптотически приближается к скорости света
    \item \textbf{Интеграция с Динамической Теорией Электромагнитного Поля Максвелла} (ДТПМ) С.А. Цикры
\end{enumerate}

%================================================================================
\section{Постановка задачи}
%================================================================================

\subsection{Исходные предположения}

Рассматривается модель, в которой:

\begin{enumerate}
    \item Электрический заряд $e$ обладает скалярным потенциалом $\phi = e/r$
    \item Поле скалярного потенциала вращается вместе с зарядом с угловой скоростью $\omega$
    \item Векторный потенциал магнитного поля определяется соотношением:
    \begin{equation}
        \mathbf{A} = \frac{\mathbf{v}}{c} \phi
    \end{equation}
    где $\mathbf{v}$ --- линейная скорость вращения поля скалярного потенциала
\end{enumerate}

\subsection{Цели исследования}

\begin{itemize}
    \item Вычислить момент инерции (угловой момент) вращающегося поля
    \item Вычислить магнитный момент системы
    \item Определить условия получения правильного гиромагнитного отношения $g = 2$
    \item Установить физический смысл радиуса перехода $R_{tr}$
    \item Проанализировать возможность интегрирования до бесконечности
    \item Интегрировать полученную модель с теорией С.А. Цикры
\end{itemize}

%================================================================================
\section{Теоретические основы}
%================================================================================

\subsection{Импульс поля по Рорлиху}

В отличие от традиционного подхода с использованием вектора Пойнтинга, в работе применяется формула импульса поля, предложенная Ф. Рорлихом, которая обеспечивает правильную релятивистскую связь между энергией и импульсом:

\begin{equation}
    \mathbf{P} = \frac{E_{field}}{c^2} \mathbf{v}
\end{equation}

Это позволяет избежать классической проблемы фактора $4/3$, когда вычисление импульса через вектор Пойнтинга даёт $\mathbf{P} = \frac{4}{3} \frac{E_{field}}{c^2} \mathbf{v}$.

\subsection{Плотность энергии электрического поля}

Для точечного заряда (или сферически симметричного распределения):

\begin{equation}
    u_E = \frac{E^2}{8\pi} = \frac{e^2}{8\pi r^4}
\end{equation}

Эффективная плотность массы поля:

\begin{equation}
    \rho_m = \frac{u_E}{c^2} = \frac{e^2}{8\pi c^2 r^4}
\end{equation}

\subsection{Характерные масштабы}

\begin{table}[h]
\centering
\caption{Характерные масштабы электрона}
\label{tab:scales}
\begin{tabular}{lcc}
\toprule
\textbf{Масштаб} & \textbf{Формула} & \textbf{Значение} \\
\midrule
Классический радиус & $r_e = \dfrac{e^2}{m c^2}$ & $2.82 \times 10^{-15}$ м \\
Комптоновская длина & $\lambda_c = \dfrac{\hbar}{m c}$ & $3.86 \times 10^{-13}$ м \\
Боровский радиус & $a_0 = \dfrac{\hbar^2}{m e^2}$ & $5.29 \times 10^{-11}$ м \\
\bottomrule
\end{tabular}
\end{table}

Соотношение между классическим радиусом и комптоновской длиной:

\begin{equation}
    \frac{r_e}{\lambda_c} = \alpha \approx \frac{1}{137}
\end{equation}

где $\alpha$ --- постоянная тонкой структуры.

%================================================================================
\section{Расчёт момента импульса поля}
%================================================================================

\subsection{Общая формула}

Момент импульса (угловой момент) вращающегося поля вычисляется как:

\begin{equation}
    \mathbf{L} = \int [\mathbf{r} \times \mathbf{g}] \, dV
\end{equation}

где $\mathbf{g}$ --- плотность импульса поля:

\begin{equation}
    \mathbf{g} = \rho_m \mathbf{v} = \frac{u_E}{c^2} \mathbf{v}
\end{equation}

\subsection{Интегрирование в сферических координатах}

В сферических координатах, учитывая, что $\mathbf{v}$ направлена по азимуту, а вклад в проекцию на ось вращения даёт $\sin^2\theta$:

\begin{equation}
    L = \int_{r_0}^{\infty} \int_{0}^{\pi} \int_{0}^{2\pi} r \cdot v(r) \sin^2\theta \cdot \frac{e^2}{8\pi c^2 r^4} \cdot r^2 \sin\theta \, dr \, d\theta \, d\phi
\end{equation}

Интеграл по углам даёт $\frac{8\pi}{3}$. Остаётся радиальный интеграл:

\begin{equation}
    L = \frac{e^2}{3 c^2} \int_{r_0}^{\infty} \frac{v(r)}{r} \, dr
\end{equation}

\subsection{Проблема расходимости}

Если $v(r) \to c$ при $r \to \infty$, то подынтегральное выражение ведёт себя как $\frac{c}{r}$, и интеграл

\begin{equation}
    \int^{\infty} \frac{c}{r} \, dr
\end{equation}

расходится \textbf{логарифмически}.

\textbf{Физический смысл:} Если всё электрическое поле электрона (простирающееся до бесконечности) вращается со скоростью света, суммарный момент импульса будет бесконечным, что противоречит конечности спина электрона ($\hbar/2$).

\subsection{Решение: введение радиуса перехода}

Вводим физическое обрезание на масштабе $R_{tr}$. Приближённо, если в основной области $v(r) \approx \omega r$:

\begin{equation}
    L \approx \frac{e^2}{3 c^2} \int_{r_0}^{R_{tr}} \omega \, dr = \frac{e^2 \omega}{3 c^2} (R_{tr} - r_0)
\end{equation}

При $R_{tr} \gg r_0$:

\begin{equation}
    L \approx \frac{e^2 \omega R_{tr}}{3 c^2}
\end{equation}

%================================================================================
\section{Расчёт магнитного момента}
%================================================================================

\subsection{Векторный потенциал}

Согласно исходному соотношению:

\begin{equation}
    \mathbf{A} = \frac{\mathbf{v}}{c} \phi = \frac{v(r)}{c} \frac{e}{r} \hat{\phi}
\end{equation}

\subsection{Оценка магнитного момента}

Магнитный момент оценивается через эффективный ток вращения заряда. Для сферически симметричного распределения, вращающегося с профилем $v(r)$:

\begin{equation}
    \mu \approx \frac{e}{2c} R_{eff} v(R_{eff})
\end{equation}

где $R_{eff}$ --- эффективный радиус вращения тока.

Поскольку скорость насыщается до $c$ на радиусе $R_{tr}$, основной вклад даёт область $R_{tr}$:

\begin{equation}
    \mu \approx \frac{e}{2c} R_{tr} c = \frac{e R_{tr}}{2}
\end{equation}

\subsection{Асимптотика магнитного поля}

На больших расстояниях ($r \gg R_{tr}$) магнитное поле должно иметь дипольную структуру:

\begin{equation}
    \mathbf{B} \sim \frac{1}{r^3}
\end{equation}

Это достигается автоматически при учёте убывания скалярного потенциала $\phi \sim 1/r$ и ограничении области вращения радиусом $R_{tr}$.

%================================================================================
\section{Гиромагнитное отношение}
%================================================================================

\subsection{Вычисление отношения}

\begin{equation}
    \frac{\mu}{L} = \frac{ \dfrac{e R_{tr}}{2} }{ \dfrac{e^2 \omega R_{tr}}{3 c^2} } = \frac{3 c^2}{2 e \omega}
\end{equation}

\subsection{Условие получения $g = 2$}

Для электрона с $g$-фактором $g \approx 2$:

\begin{equation}
    \frac{\mu}{S} = \frac{e}{mc}
\end{equation}

где $S = L$ --- момент импульса (спин).

Приравниваем:

\begin{equation}
    \frac{3 c^2}{2 e \omega} = \frac{e}{mc}
\end{equation}

Отсюда требуемая угловая скорость:

\begin{equation}
    \omega = \frac{3 m c^3}{2 e^2}
\end{equation}

\subsection{Выражение через классический радиус}

Используя определение $r_e = \dfrac{e^2}{m c^2}$:

\begin{equation}
    \frac{m c^3}{e^2} = \frac{c}{r_e}
\end{equation}

Тогда:

\begin{equation}
    \omega \approx \frac{c}{r_e}
\end{equation}

Это означает, что экваториальная скорость на классическом радиусе:

\begin{equation}
    v(r_e) = \omega r_e \approx c
\end{equation}

\textbf{Проблема:} Это противоречит условию плавного перехода! Мы хотели, чтобы $v \to c$ только на $R_{tr}$, а не на $r_e$.

\subsection{Корректное решение}

Для получения $g=2$ без сверхсветовых скоростей на малых радиусах необходимо:

\begin{enumerate}
    \item \textbf{Радиус перехода} должен совпадать с комптоновской длиной волны:
    \begin{equation}
        R_{tr} = \lambda_c = \frac{\hbar}{m c}
    \end{equation}

    \item \textbf{Угловая скорость} должна быть:
    \begin{equation}
        \omega = \frac{c}{R_{tr}} = \frac{m c^2}{\hbar}
    \end{equation}

    \item \textbf{Скорость на классическом радиусе}:
    \begin{equation}
        v(r_e) = \omega r_e = \frac{c}{\lambda_c} r_e = c \cdot \frac{r_e}{\lambda_c} = c \cdot \alpha \approx \frac{c}{137}
    \end{equation}
\end{enumerate}

Это \textbf{нерелятивистская скорость}, что физически корректно!

%================================================================================
\section{Проблема интегрирования до бесконечности}
%================================================================================

\subsection{Математическая расходимость}

При интегрировании до бесконечности с условием $v(r) \to c$:

\begin{equation}
    L \propto \int_{r_0}^{\infty} \frac{v(r)}{r} \, dr \sim \int_{r_0}^{\infty} \frac{c}{r} \, dr \to \infty
\end{equation}

\textbf{Логарифмическая расходимость} делает момент импульса бесконечным.

\subsection{Физическая несостоятельность}

Даже с использованием формулы Рорлиха модель \textbf{не работает} при интегрировании до бесконечности по следующим причинам:

\begin{enumerate}
    \item \textbf{Бесконечный спин}: $L \to \infty$, тогда как экспериментальный спин электрона конечен ($\hbar/2$)

    \item \textbf{Неправильная асимптотика магнитного поля}:
    \begin{itemize}
        \item При $v \to c$ и $\phi \sim 1/r$ получаем $\mathbf{A} \sim 1/r$
        \item Это соответствует \textbf{магнитному монополю}, а не диполю
        \item Для диполя должно быть $\mathbf{A} \sim 1/r^2$ и $\mathbf{B} \sim 1/r^3$
    \end{itemize}

    \item \textbf{Гиромагнитное отношение}: При $L \to \infty$ получаем $\mu/L \to 0$, что не соответствует $g=2$
\end{enumerate}

\subsection{Вывод}

\textbf{Интегрирование до бесконечности невозможно.} Необходим физический радиус обрезания $R_{tr}$, за пределами которого:
\begin{itemize}
    \item Поле перестаёт вращаться жёстко
    \item Скорость вращения падает
    \item Поле становится статическим
\end{itemize}

%================================================================================
\section{Радиус перехода и комптоновская длина волны}
%================================================================================

\subsection{Физическое обоснование $R_{tr} = \lambda_c$}

Комптоновская длина волны $\lambda_c = \hbar/(mc)$ выступает как естественный масштаб по следующим причинам:

\begin{enumerate}
    \item \textbf{Квантовая локализация}: Электрон не может быть локализован лучше, чем $\lambda_c$, из-за квантовых флуктуаций вакуума

    \item \textbf{Рождение пар}: На расстояниях порядка $\lambda_c$ становятся существенными процессы рождения виртуальных электрон-позитронных пар

    \item \textbf{Граница классического описания}: За пределами $\lambda_c$ классическое описание вращающегося поля теряет смысл

    \item \textbf{Частота зиттербевегунга}: Угловая скорость $\omega = c/\lambda_c = mc^2/\hbar$ соответствует частоте дрожания электрона
\end{enumerate}

%================================================================================
\section{Профиль скорости вращения поля}
%================================================================================

\subsection{Требования к профилю}

Функция $v(r)$ должна удовлетворять условиям:

\begin{enumerate}
    \item $v(r) \approx \omega r$ при $r \ll R_{tr}$ (классическое вращение)
    \item $v(r) \to c$ при $r \gg R_{tr}$ (релятивистское насыщение)
    \item Плавный переход без разрывов
    \item Обеспечение сходимости интегралов
\end{enumerate}

\subsection{Рекомендуемые формы}

\textbf{Вариант 1: Гиперболический тангенс}
\begin{equation}
    v(r) = c \cdot \tanh\left(\frac{r}{\lambda_c}\right)
\end{equation}

\textbf{Вариант 2: Алгебраическая функция}
\begin{equation}
    v(r) = c \cdot \frac{r}{\sqrt{r^2 + \lambda_c^2}}
\end{equation}

\subsection{Асимптотическое поведение}

Для варианта 2:

\begin{itemize}
    \item При $r \ll \lambda_c$:
    \begin{equation}
        v(r) \approx c \cdot \frac{r}{\lambda_c} = \omega r, \quad \omega = \frac{c}{\lambda_c}
    \end{equation}

    \item При $r \gg \lambda_c$:
    \begin{equation}
        v(r) \approx c \left(1 - \frac{\lambda_c^2}{2r^2}\right) \to c
    \end{equation}

    \item При $r = \lambda_c$:
    \begin{equation}
        v(\lambda_c) = \frac{c}{\sqrt{2}} \approx 0.707c
    \end{equation}
\end{itemize}

%================================================================================
\section{Оценка линейных скоростей}
%================================================================================

\begin{table}[h]
\centering
\caption{Линейные скорости в характерных точках}
\label{tab:velocities}
\begin{tabular}{cccccc}
\toprule
\textbf{Радиус} & $\boldsymbol{r/\lambda_c}$ & $\boldsymbol{v/c}$ & $\boldsymbol{v}$, м/с & $\boldsymbol{\gamma}$ \\
\midrule
$r_e$ (классический) & 0.0073 & 0.0073 & $2.19 \times 10^6$ & 1.0000 \\
$0.1\lambda_c$ & 0.10 & 0.0995 & $2.98 \times 10^7$ & 1.0050 \\
$0.5\lambda_c$ & 0.50 & 0.447 & $1.34 \times 10^8$ & 1.1180 \\
$\lambda_c$ (переход) & 1.00 & 0.707 & $2.12 \times 10^8$ & 1.4142 \\
$2\lambda_c$ & 2.00 & 0.894 & $2.68 \times 10^8$ & 2.2361 \\
$10\lambda_c$ & 10.0 & 0.995 & $2.98 \times 10^8$ & 10.05 \\
\bottomrule
\end{tabular}
\end{table}

\subsection{Ключевые выводы}

\begin{enumerate}
    \item \textbf{На классическом радиусе $r_e$}: Скорость нерелятивистская: $v \approx \alpha c \approx c/137$

    \item \textbf{Начало релятивистского режима}: При $r \gtrsim 0.5\lambda_c$ скорость превышает $0.4c$

    \item \textbf{Область перехода}: На $r = \lambda_c$ скорость составляет $0.707c$

    \item \textbf{Асимптотическое приближение к $c$}: Даже на $10\lambda_c$ скорость составляет 99.5\% от $c$
\end{enumerate}

%================================================================================
\section{Распределение магнитного поля}
%================================================================================

\subsection{Аналитическое выражение}

Для профиля скорости $v(r) = c \cdot \dfrac{r}{\sqrt{r^2 + \lambda_c^2}}$ и векторного потенциала $\mathbf{A} = \dfrac{\mathbf{v}}{c}\phi$:

На экваторе ($\theta = \pi/2$) магнитная индукция:

\begin{equation}
    B_\theta(r) = \frac{e \cdot \lambda_c^2}{(r^2 + \lambda_c^2)^{3/2}}
\end{equation}

\subsection{Асимптотическое поведение}

\begin{table}[h]
\centering
\caption{Поведение магнитного поля в разных областях}
\label{tab:field}
\begin{tabular}{lc}
\toprule
\textbf{Область} & \textbf{Поведение поля} \\
\midrule
$r \ll \lambda_c$ & $B \approx \text{const} = \dfrac{e}{\lambda_c^2}$ \\
$r \sim \lambda_c$ & Плавный переход \\
$r \gg \lambda_c$ & $B \propto \dfrac{1}{r^3}$ (дипольное) \\
\bottomrule
\end{tabular}
\end{table}

\subsection{Проверка дипольной асимптотики}

При $r \gg \lambda_c$:

\begin{equation}
    B_\theta(r) \approx \frac{e \lambda_c^2}{r^3} = \frac{\mu}{r^3}
\end{equation}

где $\mu = e\lambda_c$ --- характерный магнитный момент (с точностью до числового коэффициента).

\textbf{Это соответствует эксперименту!} Поле электрона на больших расстояниях --- строго дипольное.

%================================================================================
\section{Анализ теории С.А. Цикры}
%================================================================================

\subsection{Основные положения ДТПМ}

С.А. Цикра в статье ``Спин-поле-Макс'' развивает \textbf{Динамическую Теорию Электромагнитного Поля Максвелла} (ДТПМ), основанную на оригинальных работах Максвелла:

\begin{enumerate}
    \item \textbf{Векторный потенциал} $\mathbf{A}$ является импульсом электрического поля:
    \begin{equation}
        \mathbf{A} = \frac{\mathbf{v}}{c}\phi
    \end{equation}

    \item \textbf{Поле движется вместе с источником}, включая вращательное движение

    \item \textbf{Уравнение} $\mathbf{B} = \nabla \times \mathbf{A}$ показывает, что в области вокруг постоянного магнита поле потенциальное ($\nabla \times \mathbf{B} = 0$)
\end{enumerate}

\subsection{Критика традиционной теории}

Цикра указывает на фундаментальную проблему модели контурных токов:

\begin{itemize}
    \item Поле контура тока имеет \textbf{замкнутые силовые линии}, проходящие через периферию и внутренний керн в \textbf{противоположных направлениях}
    \item В среднем по объему они компенсируют друг друга
    \item Это отличается от поля электрического диполя, имеющего \textbf{одно общее направление} везде
\end{itemize}

\subsection{Поле вращающегося заряда по Цикре}

\textbf{Векторный потенциал} в плоскости вращения:
\begin{equation}
    A_\phi = \frac{q\omega}{c} = \text{const}
\end{equation}

\textbf{Ключевая особенность:} Потенциал \textbf{не убывает} с удалением от заряда!

\textbf{Магнитная индукция}:
\begin{equation}
    B_z \sim \frac{q\omega}{cr} \propto \frac{1}{r}
\end{equation}

Это отличается от дипольного поля ($\propto 1/r^3$).

\subsection{Принципиальные отличия}

\begin{table}[h]
\centering
\caption{Сравнение поля контура тока и вращающегося заряда}
\label{tab:comparison}
\begin{tabular}{lc}
\toprule
\textbf{Характеристика} & \textbf{Вращающийся заряд (Цикра)} \\
\midrule
Силовые линии B & Приходят из бесконечности и уходят в бесконечность \\
Направление & Нет обратного керна \\
Убывание B & $\propto 1/r$ \\
Потенциальность & $\nabla \times \mathbf{B} \neq 0$ (непотенциальное) \\
\bottomrule
\end{tabular}
\end{table}

\subsection{Непотенциальность поля}

\textbf{Важнейший результат Цикры:} Поле вращающегося заряда:
\begin{itemize}
    \item \textbf{Непотенциальное} ($\nabla \times \mathbf{B} \neq 0$)
    \item \textbf{Стационарное} (неизменное во времени)
\end{itemize}

Это необычно для традиционной электродинамики, где непотенциальные поля $\mathbf{B}$ могут быть только нестационарными.

\subsection{Физические следствия}

Для частицы с магнитным моментом $\boldsymbol{\mu}$ в непотенциальном поле:

\begin{equation}
    \mathbf{F} = \nabla(\boldsymbol{\mu} \cdot \mathbf{B}) + \text{дополнительный член}
\end{equation}

Второй член отсутствует для потенциальных полей, но важен для непотенциальных полей вращающихся зарядов. Это может объяснять:
\begin{itemize}
    \item Прецессию момента
    \item Осцилляции частицы
    \item \textbf{Спиновые волны} в твердых магнетиках (ранее считавшиеся чисто квантовым явлением)
\end{itemize}

\subsection{Проблема теории Цикры}

\textbf{Расходимость на бесконечности}:
\begin{itemize}
    \item Векторный потенциал не убывает: $A \sim \text{const}$
    \item Магнитное поле убывает как $1/r$, а не $1/r^3$
    \item Интегралы момента импульса и энергии расходятся
    \item Невозможно получить конечное гиромагнитное отношение
\end{itemize}

%================================================================================
\section{Интеграция моделей}
%================================================================================

\subsection{Сравнительный анализ}

\begin{table}[h]
\centering
\caption{Сравнение моделей}
\label{tab:models}
\begin{tabular}{lccc}
\toprule
\textbf{Аспект} & \textbf{Цикра} & \textbf{Наша модель} & \textbf{Объединение} \\
\midrule
Что вращается & Поле заряда & Скалярный $\phi$ & $\mathbf{A} = \frac{\mathbf{v}}{c}\phi$ \\
Профиль скорости & $v = \omega r$ & Плавный к $c$ & Насыщение на $\lambda_c$ \\
Асимптотика $\mathbf{A}$ & $\sim \text{const}$ & $\sim 1/r$ & $\sim 1/r$ \\
Асимптотика $\mathbf{B}$ & $\sim 1/r$ & $\sim 1/r^3$ & $\sim 1/r^3$ \\
Обрезание & Нет & $R_{tr} = \lambda_c$ & $R_{tr} = \lambda_c$ \\
Гиромагнитное отнош. & Не вычислено & $g = 2$ & $g = 2$ \\
\bottomrule
\end{tabular}
\end{table}

\subsection{Схема объединённой модели}

\textbf{Шаг 1: Базовое уравнение связи}
\begin{equation}
    \mathbf{A}(\mathbf{r}) = \frac{\mathbf{v}(\mathbf{r})}{c} \, \phi(\mathbf{r}), \quad \phi(\mathbf{r}) = \frac{q}{r}
\end{equation}

\textbf{Шаг 2: Профиль скорости вращения поля}
\begin{equation}
    v(r) = c \cdot \tanh\left(\frac{r}{\lambda_c}\right) \quad \text{или} \quad v(r) = c \cdot \frac{r}{\sqrt{r^2 + \lambda_c^2}}
\end{equation}

\textbf{Шаг 3: Вычисление магнитного поля}
\begin{equation}
    \mathbf{B} = \nabla \times \mathbf{A} = \nabla \times \left( \frac{\mathbf{v}}{c} \phi \right)
\end{equation}

\textbf{Шаг 4: Момент импульса поля (по Рорлиху)}
\begin{equation}
    \mathbf{L} = \int \left[ \mathbf{r} \times \left( \frac{u_E}{c^2} \mathbf{v} \right) \right] dV, \quad u_E = \frac{E^2}{8\pi}
\end{equation}

\textbf{Шаг 5: Магнитный момент}
\begin{equation}
    \boldsymbol{\mu} = \frac{1}{2c} \int [\mathbf{r} \times \mathbf{j}_{eff}] dV, \quad \mathbf{j}_{eff} = \rho \mathbf{v}
\end{equation}

\textbf{Шаг 6: Гиромагнитное отношение}
\begin{equation}
    \frac{\mu}{L} = \frac{q}{m c} \quad \Rightarrow \quad g = 2
\end{equation}

%================================================================================
\section{Основные результаты и выводы}
%================================================================================

\subsection{Главные результаты}

\begin{enumerate}
    \item \textbf{Получено правильное гиромагнитное отношение $g = 2$} в рамках классической модели при условии:
    \begin{itemize}
        \item Радиус перехода $R_{tr} = \lambda_c$ (комптоновская длина волны)
        \item Угловая скорость $\omega = c/\lambda_c$
        \item Плавный профиль скорости с насыщением к $c$
    \end{itemize}

    \item \textbf{Доказана невозможность интегрирования до бесконечности}:
    \begin{itemize}
        \item Логарифмическая расходимость момента импульса
        \item Неправильная асимптотика магнитного поля ($1/r$ вместо $1/r^3$)
        \item Невозможность получения конечного $g$-фактора
    \end{itemize}

    \item \textbf{Установлен физический смысл радиуса перехода}:
    \begin{itemize}
        \item $\lambda_c$ --- естественная граница применимости классического описания
        \item На масштабах $r > \lambda_c$ существенны квантовые эффекты
        \item На масштабах $r < \lambda_c$ доминирует классическое вращение поля
    \end{itemize}

    \item \textbf{Получены реалистичные оценки скоростей}:
    \begin{itemize}
        \item На классическом радиусе: $v \approx c/137$ (нерелятивистская)
        \item На комптоновском радиусе: $v \approx 0.707c$ (релятивистская)
        \item Асимптотическое приближение к $c$ без достижения
    \end{itemize}

    \item \textbf{Воспроизведена правильная структура магнитного поля}:
    \begin{itemize}
        \item Дипольная асимптотика $\mathbf{B} \sim 1/r^3$ на больших расстояниях
        \item Сложная непотенциальная структура в ближней зоне
        \item Отсутствие обратного керна (в соответствии с теорией Цикры)
    \end{itemize}
\end{enumerate}

\subsection{Интеграция с теорией Цикры}

\begin{enumerate}
    \item \textbf{Теория Цикры} даёт важную физическую идею: вращение поля вместе с зарядом порождает магнитное поле с уникальной непотенциальной структурой

    \item \textbf{Без обрезания} на комптоновском масштабе теория Цикры приводит к расходимостям и не может дать конечное гиромагнитное отношение

    \item \textbf{Объединённая модель} сохраняет преимущества обоих подходов:
    \begin{itemize}
        \item Непотенциальность поля в ближней зоне (от Цикры)
        \item Конечность физических величин и правильное $g=2$ (от нашей модели)
        \item Дипольная асимптотика на больших расстояниях
    \end{itemize}
\end{enumerate}

\subsection{Физические следствия}

\begin{enumerate}
    \item \textbf{Спин электрона} --- это момент импульса вращающегося электромагнитного поля, локализованный в области $\sim \lambda_c$

    \item \textbf{Магнитный момент} формируется тем же вращением, но с правильной дипольной асимптотикой

    \item \textbf{Непотенциальность поля} в ближней зоне объясняет квантоподобные эффекты:
    \begin{itemize}
        \item Спиновые волны в магнетиках
        \item Обменное взаимодействие
        \item Прецессию спина
    \end{itemize}
    \textbf{без привлечения постулатов квантовой механики}

    \item \textbf{Комптоновская длина волны} выступает как естественная граница между классическим и квантовым описаниями
\end{enumerate}

\subsection{Перспективы развития}

\begin{enumerate}
    \item Количественное описание спиновых волн в рамках классической теории
    \item Исследование обменного взаимодействия как следствия непотенциальности поля
    \item Расчёт тонкой структуры уровней атома водорода с учётом классического спина
    \item Связь с квантовой электродинамикой: построение классического предела КЭД
    \item Экспериментальная проверка предсказанных отклонений на масштабах порядка $\lambda_c$
\end{enumerate}

%================================================================================
\section{Приложения}
%================================================================================

\subsection{Основные константы}

\begin{table}[h]
\centering
\caption{Физические константы (СИ)}
\label{tab:constants}
\begin{tabular}{lc}
\toprule
\textbf{Константа} & \textbf{Значение} \\
\midrule
Скорость света $c$ & $2.998 \times 10^8$ м/с \\
Заряд электрона $e$ & $1.602 \times 10^{-19}$ Кл \\
Масса электрона $m$ & $9.109 \times 10^{-31}$ кг \\
Постоянная Планка $\hbar$ & $1.055 \times 10^{-34}$ Дж·с \\
Постоянная тонкой структуры $\alpha$ & $1/137.036$ \\
\bottomrule
\end{tabular}
\end{table}

\subsection{Основные уравнения модели}

\begin{enumerate}
    \item \textbf{Векторный потенциал}:
    \begin{equation}
        \mathbf{A} = \frac{\mathbf{v}}{c} \phi, \quad \phi = \frac{e}{r}
    \end{equation}

    \item \textbf{Профиль скорости}:
    \begin{equation}
        v(r) = c \cdot \frac{r}{\sqrt{r^2 + \lambda_c^2}}
    \end{equation}

    \item \textbf{Магнитное поле} (на экваторе):
    \begin{equation}
        B_\theta(r) = \frac{e \lambda_c^2}{(r^2 + \lambda_c^2)^{3/2}}
    \end{equation}

    \item \textbf{Момент импульса}:
    \begin{equation}
        L = \frac{e^2}{3c^2} \int_{r_e}^{\lambda_c} \frac{v(r)}{r} dr
    \end{equation}

    \item \textbf{Магнитный момент}:
    \begin{equation}
        \mu \approx \frac{e \lambda_c}{2}
    \end{equation}

    \item \textbf{Гиромагнитное отношение}:
    \begin{equation}
        \frac{\mu}{L} = \frac{e}{mc} \quad \Rightarrow \quad g = 2
    \end{equation}
\end{enumerate}

%================================================================================
\section*{Заключение}
%================================================================================

Разработанная классическая модель спина электрона на основе вращающегося скалярного потенциала:

\begin{itemize}
    \item[\checkmark] \textbf{Воспроизводит правильное гиромагнитное отношение} $g = 2$
    \item[\checkmark] \textbf{Избегает сверхсветовых скоростей} на классическом радиусе
    \item[\checkmark] \textbf{Обеспечивает правильную дипольную асимптотику} магнитного поля
    \item[\checkmark] \textbf{Объясняет квантоподобные эффекты} в рамках классической теории
    \item[\checkmark] \textbf{Интегрируется с ДТПМ} С.А. Цикры
\end{itemize}

\textbf{Ключевой результат:} Комптоновская длина волны $\lambda_c$ выступает как физически обоснованный радиус перехода, за пределами которого классическое описание вращающегося поля теряет применимость.

\textbf{Перспектива:} Данная модель открывает путь к построению непротиворечивого классического предела квантовой электродинамики, где спин и магнитный момент электрона получают естественное объяснение через вращение электромагнитного поля.

%================================================================================
\begin{thebibliography}{99}
%================================================================================

\bibitem{landau}
Л.Д. Ландау, Е.М. Лифшиц. \textit{Теоретическая физика. Т.8. Электродинамика сплошных сред.} М.: Наука, 1982.

\bibitem{maxwell1}
Д.К. Максвелл. \textit{Избранные труды по теории электромагнитного поля.} М.: Гостехиздат, 1952.

\bibitem{maxwell2}
Д.К. Максвелл. \textit{Трактат об электричестве и магнетизме. Т.2.} М.: Наука, 1989.

\bibitem{tsykra}
С.А. Цикра. \textit{Спин-поле-Макс.} Донецк, Украина, 2026.

\bibitem{feynman}
Р. Фейнман, Р. Лейтон, М. Сэндс. \textit{Фейнмановские лекции по физике. Т.5.} М.: Наука, 1966.

\bibitem{akhiezer}
А.И. Ахиезер, В.Г. Барьяхтар, С.В. Пелетминский. \textit{Спиновые волны.} М.: Наука, 1965.

\bibitem{rohrlich}
F. Rohrlich. \textit{Classical Charged Particles.} Addison-Wesley, 1965.

\end{thebibliography}

\end{document}